\chapter{Całka Stratonovicha i jej związek z całką Itô}
\label{ch:stratonovich}

\section{Dwie reguły sumowania, dwa rachunki}
W klasycznym rachunku różniczkowym funkcja złożona
spełnia regułę łańcuchową bez dodatkowych wyrazów.
W rachunku Itô pojawia się poprawka zależna od wariacji
kwadratowej. Można jednak zdefiniować inną całkę,
która w przypadku ciągłych procesów Itô zachowuje
zwykłą regułę łańcuchową. Jest nią \emph{całka
Stratonovicha}.

Trzeba tu odróżnić dwa pytania. Pierwsze dotyczy
\emph{istnienia granicy} sum, w których współczynnik
wybieramy na środku przedziału. Drugie dotyczy tego,
\emph{co dokładnie rozumiemy przez środek}: w poniższej
konstrukcji użyjemy średniej arytmetycznej wartości
procesu podcałkowego na obu końcach przedziału.
Takie sumy są symetryczne, lecz w przeciwieństwie do
sum Itô korzystają z przyszłej wartości integrandu.
Nie jest to strategia przewidywalna w rozumieniu
rozdziału \ref{ch:ito-construction}.

\begin{przyklad}[Najprostsza różnica]
Dla procesu Wienera $W_0=0$ rozważmy
\[
S_\pi^{\mathrm L}=\sum_kW_{t_k}\Delta_kW,
\qquad
S_\pi^{\mathrm S}=
\frac12\sum_k(W_{t_k}+W_{t_{k+1}})\Delta_kW.
\]
Drugie wyrażenie jest dokładnie teleskopowe:
$S_\pi^{\mathrm S}=\frac12W_T^2$ dla każdego podziału.
Natomiast $S_\pi^{\mathrm L}\to\frac12(W_T^2-T)$
w $L^2$. Zatem odpowiednia poprawka wynosi $T/2$.
\emph{To nie wybór dowolnej estetycznej konwencji}:
obydwa rachunki odpowiadają różnym konstrukcjom
matematycznym i różnym sposobom aproksymowania hałasu.
\end{przyklad}

\begin{center}
\begin{tikzpicture}[>=Latex,x=1.0cm,y=.85cm]
\draw[->] (0,0)--(8.5,0) node[right] {$t$};
\draw[->] (0,-.35)--(0,2.5) node[above] {$W_t$};
\draw[thick,blue!70!black]
 (0,0)--(.6,.3)--(1.2,1.15)--(1.8,.9)--(2.4,1.6)
 --(3,1.0)--(3.6,.65)--(4.2,1.35)--(4.8,1.1)
 --(5.4,1.9)--(6,1.4)--(6.6,2.1)--(7.2,1.5)--(7.8,1.85);
\foreach \x in {0,1.8,3.6,5.4,7.8}
 {\draw[dashed,gray!70] (\x,0)--(\x,2.25);}
\fill[blue!65!black] (1.8,.9) circle (2pt);
\fill[blue!65!black] (3.6,.65) circle (2pt);
\draw[<->,orange!75!black,thick] (1.8,2.3)--(3.6,2.3)
 node[midway,above,scale=.85] {jeden przedział};
\node[scale=.81,anchor=north] at (1.8,-.12) {$t_k$};
\node[scale=.81,anchor=north] at (3.6,-.12) {$t_{k+1}$};
\node[scale=.78,align=left,anchor=west] at (5.0,2.42)
 {Itô: lewy koniec\\Stratonovich: średnia końców};
\end{tikzpicture}
\par\smallskip\small Porównanie reguł próbkowania na tym samym odcinku trajektorii.
\end{center}

\section{Kowariacja i zbieżność sum symetrycznych}
\begin{definicja}[Kowariacja ciągłych procesów Itô]
Niech $H$ i $X$ będą ciągłymi procesami Itô,
$H=H_0+\int_0^.\alpha_s\,\dd s+
\int_0^.\beta_s\,\dd W_s$ oraz
$X=X_0+\int_0^.b_s\,\dd s+
\int_0^.\sigma_s\,\dd W_s$,
z lokalnymi warunkami całkowalności jak w
rozdziale \ref{ch:ito-formula}. Definiujemy
\[
[H,X]_t:=\int_0^t\beta_s\sigma_s\,\dd s.
\tag{4.1}\label{eq:HX-bracket}
\]
Jest to dobrze określony ciągły proces skończonej
wariacji, ponieważ
$\int_0^T|\beta_s\sigma_s|\,\dd s
\le(\int_0^T\beta_s^2\,\dd s)^{1/2}
(\int_0^T\sigma_s^2\,\dd s)^{1/2}<\infty$
prawie na pewno.
\end{definicja}

\begin{twierdzenie}[Zbieżność symetrycznych sum]
\label{th:strat-sums}
Niech $X,H$ będą ciągłymi procesami Itô, jak w
poprzedniej definicji. Niech $\pi_n$ będą
\emph{deterministycznymi} podziałami $[0,T]$
o średnicy dążącej do zera. Dla $t\in[0,T]$
niech
\[
S_t^{(n)}:=\frac12\sum_k
(H_{t_k}+H_{t_{k+1}})
\bigl(X_{t_{k+1}\wedge t}-X_{t_k\wedge t}\bigr)
\ind{\{t_k<t\}},
\]
przy czym w składniku przeciętym przez $t$
zastępujemy $H_{t_{k+1}}$ wartością $H_t$.
Wtedy
\[
S_t^{(n)}\ \xrightarrow[n\to\infty]{\mathrm{ucp}}
\int_0^tH_s\,\dd X_s+\frac12[H,X]_t,
\tag{4.2}\label{eq:strat-sum}
\]
gdzie $\int H\,\dd X:=\int Hb\,\dd s+
\int H\sigma\,\dd W$ oznacza całkę Itô.
\end{twierdzenie}
\begin{proof}
Dla $t=T$ rozkładamy sumę na część lewostronną i
część poprawkową:
\[
S_T^{(n)}=\sum_kH_{t_k}\Delta_kX
+\frac12\sum_k\Delta_kH\,\Delta_kX.
\tag{4.3}\label{eq:strat-decomp}
\]
Dla składnika pierwszego definiujemy
$H_s^{(n)}=H_{t_k}$ na $(t_k,t_{k+1}]$.
Jest to proces przewidywalny. Z ciągłości $H$ mamy
$\sup_s|H_s^{(n)}-H_s|\to0$ ścieżkowo. Część
całki względem skończonej wariacji spełnia
\[
\left|\int_0^T(H_s^{(n)}-H_s)b_s\,\dd s\right|
\le\norm{H^{(n)}-H}_\infty\int_0^T|b_s|\,\dd s
\longrightarrow0
\]
prawie na pewno. Dla części stochastycznej
lokalizujemy energię $\int_0^T\sigma_s^2\,\dd s$
i supremum $|H|$. Wtedy zbieżność zdominowana
i izometria Itô zapewniają
$\E|\int_0^T(H_s^{(n)}-H_s)\sigma_s\,\dd W_s|^2\to0$
po zatrzymaniu. Usuwając lokalizację, otrzymujemy
zbieżność w prawdopodobieństwie do $\int_0^TH_s\,\dd X_s$.

W drugim składniku \eqref{eq:strat-decomp} korzystamy
z kowariacji. Z polaryzacji i lematu
\ref{lem:ito-qv} zastosowanego do $H+X$ i $H-X$
wynika ucp-zbieżność
$\sum_{t_{k+1}\le t}\Delta_kH\Delta_kX\to[H,X]_t$.
Połączenie obydwu granic daje \eqref{eq:strat-sum}
w chwili końcowej. Dla uciętych sum wszystkie powyższe
oszacowania zachodzą także w supremum po $t$, ponieważ
całki stochastyczne kontroluje nierówność maksymalna
Dooba, a sumy kwadratów przyrostów zbiegają ucp.
Różnica spowodowana ostatnim, niepełnym odcinkiem
jest zaniedbywalna wskutek ciągłości $H,X$ i
ograniczoności odpowiednich sum w prawdopodobieństwie.
Otrzymujemy zbieżność ucp na całym $[0,T]$.
\end{proof}

\begin{definicja}[Całka Stratonovicha]
Dla ciągłych procesów Itô $H,X$ z poprzedniego
twierdzenia definiujemy
\[
\int_0^tH_s\circ\dd X_s
:=\int_0^tH_s\,\dd X_s+\frac12[H,X]_t.
\tag{4.4}\label{eq:strat-def}
\]
Symbol $\circ\dd X_s$ oznacza całkowanie
Stratonovicha, natomiast $\dd X_s$ bez kółka
oznacza całkowanie Itô. Równość
\eqref{eq:strat-sum} pokazuje, że definicja zgadza
się z granicą symetrycznych sum Riemanna.
\end{definicja}

\begin{uwaga}[Granice stosowalności]
Wzór \eqref{eq:strat-def} zdefiniowaliśmy dla
\emph{ciągłych procesów Itô}, a nie dla każdego
procesu przewidywalnego. Dla ogólnych
półmartyngałów ciągłych można go rozszerzyć przez
kowariację; dla procesów ze skokami trzeba ponadto
rozróżnić kilka nierównoważnych konwencji (np.
całkę kanoniczną Marcusa). Nie będziemy ich utożsamiać.
\end{uwaga}

\section{Reguła łańcuchowa bez członu drugiego rzędu}
\begin{twierdzenie}[Reguła łańcuchowa Stratonovicha]
\label{th:strat-chain}
Niech $X$ będzie jednowymiarowym ciągłym procesem
Itô. Dla każdej funkcji $F\in C^3(\R)$ zachodzi
prawie na pewno, dla wszystkich $t\in[0,T]$,
\[
F(X_t)-F(X_0)=\int_0^t F'(X_s)\circ\dd X_s.
\tag{4.5}\label{eq:strat-chain}
\]
Jeśli $F\in C^{1,3}([0,T]\times\R)$, a pochodna mieszana
$F_{tx}$ istnieje i jest ciągła, to
\[
F(t,X_t)-F(0,X_0)
=\int_0^tF_t(s,X_s)\,\dd s
+\int_0^tF_x(s,X_s)\circ\dd X_s.
\tag{4.6}\label{eq:strat-chain-time}
\]
\end{twierdzenie}
\begin{proof}
Ze wzoru Itô mamy
$F(X_t)-F(X_0)=\int_0^tF'(X_s)\,\dd X_s
+\frac12\int_0^tF''(X_s)\,\dd[X]_s$.
Ponadto, stosując wzór Itô do $F'(X)$
(co wyjaśnia założenie $F\in C^3$),
dostajemy, że współczynnik przy $\dd W_s$ w
rozkładzie $F'(X_s)$ jest równy
$F''(X_s)\sigma_s$. Stąd
$[F'(X),X]_t=\int_0^tF''(X_s)\sigma_s^2\,\dd s
=\int_0^tF''(X_s)\,\dd[X]_s$.
Po podstawieniu definicji \eqref{eq:strat-def}
człony drugiego rzędu znoszą się i otrzymujemy
\eqref{eq:strat-chain}. W przypadku funkcji zależnej
od czasu stosujemy tę samą argumentację do
$H_s=F_x(s,X_s)$: skończona wariacja składnika
$\int F_{xt}\,\dd s$ nie wnosi kowariacji.
\end{proof}

\begin{zadanie}[Trzy elementarne całki]
Obliczyć, bez pozostawiania symboli całek po prawej
stronie, wyrażenia
$\int_0^TW_s\circ\dd W_s$,
$\int_0^TW_s^2\circ\dd W_s$
oraz $\int_0^Te^{W_s}\circ\dd W_s$.

\emph{Rozwiązanie.} Z reguły łańcuchowej dla
$F(x)=x^2/2$, $F(x)=x^3/3$ i $F(x)=e^x$
(odpowiednio) otrzymujemy
\[
\int_0^TW_s\circ\dd W_s=\frac12W_T^2,
\quad \int_0^TW_s^2\circ\dd W_s=\frac13W_T^3,
\quad \int_0^Te^{W_s}\circ\dd W_s=e^{W_T}-1.
\]
Dla pierwszej całki porównanie z Itô daje
$\int_0^T W_s\circ\dd W_s
=\int_0^TW_s\,\dd W_s+T/2$.
Dla drugiej różnica wynosi $\int_0^TW_s\,\dd s$;
nie jest zatem stałą deterministyczną.
\end{zadanie}

\begin{twierdzenie}[Iloczyn i całkowanie przez części]
Dla dwóch ciągłych procesów Itô $X,Y$ zachodzi
\[
X_tY_t-X_0Y_0
=\int_0^tX_s\circ\dd Y_s+
\int_0^tY_s\circ\dd X_s.
\tag{4.7}\label{eq:strat-product}
\]
\end{twierdzenie}
\begin{proof}
Ze wzoru na iloczyn Itô (twierdzenie
\ref{th:ito-product}) otrzymujemy
$\dd(XY)=X\,\dd Y+Y\,\dd X+\dd[X,Y]$.
Z kolei \eqref{eq:strat-def} daje
$\int X\circ\dd Y=\int X\,\dd Y+\frac12[X,Y]$
oraz analogiczną równość po zamianie $X,Y$.
Dodanie obu wyrażeń daje dokładnie
\eqref{eq:strat-product}.
\end{proof}

\begin{lemat}[Jednoznaczność rozkładu na martyngał i skończoną wariację]
\label{lem:uniqueness-mart-fv}
Niech $M,N$ będą ciągłymi lokalnymi martyngałami zaczynającymi się
w zerze, a $A,B$ ciągłymi, adaptowanymi procesami o skończonej
wariacji na każdym przedziale zwartym, także zaczynającymi się
w zerze. Jeżeli $M_t+A_t=N_t+B_t$ dla wszystkich $t$ prawie
na pewno (z dokładnością do nierozróżnialności), to $M=N$ i
$A=B$ z tą samą dokładnością.
\end{lemat}
\begin{proof}
Proces $D=M-N=B-A$ jest ciągłym lokalnym martyngałem
i jednocześnie ma skończoną wariację. Dla dowolnej
trajektorii ciągłego procesu o skończonej wariacji na $[0,T]$
\[
\sum_{[u,v]\in\pi}(D_v-D_u)^2
\le\max_{[u,v]\in\pi}|D_v-D_u|\,\operatorname{TV}_{[0,T]}(D)
\longrightarrow0
\]
gdy średnica deterministycznego podziału dąży do zera.
Z definicji wariacji kwadratowej ciągłego lokalnego
martyngału wynika więc $[D]_t=0$ dla każdego $t$.
Zatrzymajmy $D$ tak, aby $D^\tau$ i jego kwadrat były
całkowalne. Proces $(D^\tau)^2-[D^\tau]$
jest martyngałem, zatem
$\E(D_{t\wedge\tau})^2=0$.
Otrzymujemy $D_{t\wedge\tau}=0$ prawie na pewno.
Usuwamy lokalizację i korzystamy z ciągłości trajektorii,
aby uzyskać $D\equiv0$, co dowodzi obu równości.
\end{proof}

\section{Konwersja równań Itô i Stratonovicha}
Sama funkcja $\sigma$ nie określa procesu, jeśli nie
wiemy, w jakiej konwencji jest całka. Rozważmy
formalnie równania
\[
\dd X_t=b_I(X_t)\,\dd t+\sigma(X_t)\,\dd W_t,
\qquad
\dd X_t=b_S(X_t)\,\dd t+\sigma(X_t)\circ\dd W_t.
\]
Współczynniki dryfu $b_I$ i $b_S$ nie są na ogół
identyczne.

\begin{twierdzenie}[Konwersja w jednym wymiarze]
\label{th:ito-strat-conversion}
Niech $b_S,\sigma\in C^2(\R)$ mają lokalnie ograniczone
pochodne, a $X$ będzie ciągłym procesem Itô o
wartościach rzeczywistych, dla którego
\[
X_t=X_0+\int_0^tb_S(X_s)\,\dd s
+\int_0^t\sigma(X_s)\circ\dd W_s
\tag{4.8}\label{eq:StratSDE}
\]
przy czym wszystkie całki istnieją lokalnie.
Wówczas równoważnie $X$ spełnia równanie Itô
\[
X_t=X_0+\int_0^t
\left[b_S(X_s)+\frac12\sigma(X_s)\sigma'(X_s)\right]\dd s
+\int_0^t\sigma(X_s)\,\dd W_s.
\tag{4.9}\label{eq:ItoSDE}
\]
Odwrotnie, każde rozwiązanie Itô \eqref{eq:ItoSDE}
(mające wskazaną regularność) spełnia \eqref{eq:StratSDE}.
\end{twierdzenie}
\begin{proof}
Najpierw wyeliminujmy pozorną kolistość rozumowania.
Z założenia $X$ jest procesem Itô, zatem istnieją
procesy $\widetilde b$ i $\widetilde\sigma$, dla których
$\dd X_t=\widetilde b_t\,\dd t+
\widetilde\sigma_t\,\dd W_t$. Po rozwinięciu całki
Stratonovicha z \eqref{eq:StratSDE} otrzymujemy
\[
X_t=X_0+\int_0^t b_S(X_s)\,\dd s
+\int_0^t\sigma(X_s)\,\dd W_s
+\tfrac12[\sigma(X),W]_t.
\]
Ostatni składnik ma skończoną wariację. Na mocy lematu \ref{lem:uniqueness-mart-fv}
jednoznaczność rozkładu na ciągły lokalny martyngał
i proces o skończonej wariacji daje
$\int_0^.\widetilde\sigma_s\,\dd W_s
=\int_0^.\sigma(X_s)\,\dd W_s$.
Z izometrii po lokalizacji wnioskujemy
$\widetilde\sigma_s=\sigma(X_s)$ dla
$(\lambda_1\otimes\Prob)$-prawie wszystkich $(s,\omega)$.
Teraz wzór Itô dla $\sigma(X)$ pokazuje, że jego
współczynnik dyfuzji wynosi
$\sigma'(X_s)\sigma(X_s)$; zatem
$[\sigma(X),W]_t=\int_0^t\sigma'(X_s)\sigma(X_s)\,\dd s$.
Definicja całki Stratonovicha daje więc
\[
\int_0^t\sigma(X_s)\circ\dd W_s
=\int_0^t\sigma(X_s)\,\dd W_s
+\frac12\int_0^t\sigma'(X_s)\sigma(X_s)\,\dd s.
\]
Po podstawieniu do \eqref{eq:StratSDE} otrzymujemy
\eqref{eq:ItoSDE}. W drugą stronę odejmujemy tę samą
poprawkę. Lokalizacja pozwala kontrolować współczynniki
bez zakładania ich globalnej ograniczoności.
\end{proof}

\begin{zadanie}[Geometryczny ruch Browna w obu konwencjach]
Niech $X_0=x_0>0$, $\mu\in\R$, $\gamma\in\R$.
Pokazać, że proces
$X_t=x_0\exp((\mu-\gamma^2/2)t+\gamma W_t)$
rozwiązuje równanie Itô
$\dd X_t=\mu X_t\,\dd t+\gamma X_t\,\dd W_t$
i podać równoważne równanie Stratonovicha.

\emph{Rozwiązanie.} Stosujemy wzór Itô do
$F(t,w)=x_0\exp((\mu-\gamma^2/2)t+\gamma w)$.
Ponieważ $F_t=(\mu-\gamma^2/2)F$,
$F_w=\gamma F$ i $F_{ww}=\gamma^2F$, wynika stąd
$\dd X_t=\mu X_t\,\dd t+\gamma X_t\,\dd W_t$.
Równoważny zapis Stratonovicha brzmi
\[
\dd X_t=(\mu-\tfrac12\gamma^2)X_t\,\dd t
+\gamma X_t\circ\dd W_t.
\]
Przy zerowym dryfie \emph{Stratonovicha}, czyli
$\dd X_t=\gamma X_t\circ\dd W_t$, rozwiązaniem jest
$x_0e^{\gamma W_t}$, nie
$x_0e^{\gamma W_t-\gamma^2t/2}$.
\end{zadanie}

\begin{przyklad}[Szum addytywny]
Jeśli $\sigma(x)\equiv\gamma$ jest stała, to
$\sigma\sigma'=0$. Równania
$\dd X_t=-\kappa X_t\,\dd t+\gamma\,\dd W_t$
oraz
$\dd X_t=-\kappa X_t\,\dd t+\gamma\circ\dd W_t$
są równoważne. Dotyczy to m.in. procesu
Ornsteina--Uhlenbecka z szumem addytywnym.
Różnica między konwencjami pojawia się dopiero przy
odpowiedniej zależności współczynnika dyfuzji od stanu.
\end{przyklad}

\section{Wielowymiarowa poprawka i znaczenie geometryczne}
\begin{twierdzenie}[Konwersja wielowymiarowa]
Niech $W=(W^1,\ldots,W^m)$ będzie $m$-wymiarowym
ruchem Browna, $X$ ciągłym procesem Itô w $\R^d$,
a $b_S:\R^d\to\R^d$ i
$\sigma:\R^d\to\R^{d\times m}$ będą klasy $C^2$
z lokalnie ograniczonymi pochodnymi.
Załóżmy, że $X$ spełnia w sensie całkowym
\[
\dd X_t^i=b_S^i(X_t)\,\dd t
+\sum_{k=1}^m\sigma_{ik}(X_t)\circ\dd W_t^k,
\quad i=1,\ldots,d.
\]
Wówczas $X$ spełnia równanie Itô o tym samym
współczynniku dyfuzji oraz dryfie
\[
b_I^i(x)=b_S^i(x)+\frac12\sum_{j=1}^d
\sum_{k=1}^m\partial_{x_j}\sigma_{ik}(x)\,
\sigma_{jk}(x).
\tag{4.10}\label{eq:strat-multi-correction}
\]
\end{twierdzenie}
\begin{proof}
Podobnie jak w dowodzie jednowymiarowym, zaczynamy
od dowolnego rozkładu Itô procesu $X$ o macierzy
dyfuzji $\widetilde\sigma$. Zapis całkowy równania
Stratonovicha różni się od całki Itô z macierzą
$\sigma(X)$ wyłącznie procesem o skończonej wariacji.
Lemat \ref{lem:uniqueness-mart-fv}, a następnie izometria
Itô, dają $\widetilde\sigma_s=\sigma(X_s)$ prawie
wszędzie względem $\dd s\otimes\dd\Prob$.
Dopiero teraz, dla ustalonego $i,k$, stosujemy
wielowymiarowy wzór Itô do $H_t=\sigma_{ik}(X_t)$.
Współczynnik przy $\dd W_t^\ell$ w jego części
martyngałowej jest równy
$\sum_j\partial_{x_j}\sigma_{ik}(X_t)\sigma_{j\ell}(X_t)$.
Zatem, z definicji kowariacji,
$[\sigma_{ik}(X),W^k]_t=
\int_0^t\sum_j\partial_{x_j}\sigma_{ik}(X_s)
\sigma_{jk}(X_s)\,\dd s$.
Sumujemy to po $k$ i dodajemy połowę kowariacji do
całek Itô. To jest dokładnie różnica dryfów podana
we wzorze \eqref{eq:strat-multi-correction}.
\end{proof}

\begin{przyklad}[Pole wektorowe na płaszczyźnie]
Niech $V(x,y)=(-y,x)$ i rozważmy równanie
Stratonovicha $\dd X_t=V(X_t)\circ\dd W_t$
dla $X_t=(X_t^1,X_t^2)$. Macierz pochodnej $DV$
ma postać $\left(\begin{smallmatrix}0&-1\\1&0\end{smallmatrix}\right)$,
zatem $DV(x,y)V(x,y)=(-x,-y)$.
Równoważne równanie Itô to
\[
\dd X_t=-\frac12X_t\,\dd t+V(X_t)\,\dd W_t.
\]
Z reguły łańcuchowej Stratonovicha dla $\norm{x}^2$
wynika $\dd\norm{X_t}^2=2X_t\cdot V(X_t)\circ\dd W_t=0$.
Stąd $\norm{X_t}=\norm{X_0}$ dla wszystkich $t$:
trajektoria porusza się po okręgu. W konwencji Itô
ten sam wynik wymaga poprawki od wariacji kwadratowej:
$-\norm{X_t}^2\,\dd t+\norm{V(X_t)}^2\,\dd t=0$.
To elementarny przykład geometrycznego znaczenia
rachunku Stratonovicha.
\end{przyklad}

\section{Aproksymacje gładkimi trajektoriami: elementarny przypadek Wong--Zakaia}
\begin{zadanie}[Dlaczego wygładzenie szumu prowadzi do Stratonovicha?]
Niech $W^{(n)}$ będzie liniową interpolacją
trajektorii Wienera na deterministycznym podziale
$\pi_n$ o średnicy dążącej do zera. Dla $x_0>0$
i $\gamma\in\R$ rozwiązać klasyczne równanie
Riemanna--Stieltjesa
\[
X_t^{(n)}=x_0+\gamma\int_0^t
X_s^{(n)}\,\dd W_s^{(n)}
\]
i znaleźć granicę jednostajną na $[0,T]$.

\emph{Rozwiązanie.} Każda trajektoria $W^{(n)}$
jest odcinkami liniowa, a więc ma skończoną wariację
na $[0,T]$. Klasyczna reguła łańcuchowa daje
$X_t^{(n)}=x_0\exp(\gamma W_t^{(n)})$.
Ponieważ prawie każda trajektoria $W$ jest
jednostajnie ciągła na $[0,T]$, jej interpolacje
liniowe zbiegają jednostajnie do $W$. Ciągłość
funkcji wykładniczej na zwartym zbiorze zawierającym
wartości wszystkich interpolacji zapewnia
\[
\sup_{t\le T}|X_t^{(n)}-x_0e^{\gamma W_t}|
\longrightarrow0\qquad\text{prawie na pewno}.
\]
Proces graniczny rozwiązuje
$\dd X_t=\gamma X_t\circ\dd W_t$,
a nie $\dd X_t=\gamma X_t\,\dd W_t$.
Udowodniliśmy zatem elementarną wersję zjawiska
Wong--Zakaia. Ogólne twierdzenie dla wielowymiarowych,
nieliniowych układów wymaga dodatkowej teorii i
nie wynika samo z tego przykładu.
\end{zadanie}

\section{Co otrzymaliśmy i jaką teorię rozwiniemy dalej?}
Całka Stratonovicha nie jest prostym przemianowaniem
całki Itô. Łączy się z nią inna reguła sumowania i
inne przesunięcie współczynnika dryfu. Jest wygodna
przy zmianie współrzędnych, w problemach geometrycznych
i przy aproksymowaniu ruchu Browna gładkimi ścieżkami.
Natomiast całka Itô naturalnie współpracuje z
warunkową wartością oczekiwaną, martyngałami i
rachunkiem probabilistycznym. Obydwie konwencje są
ściśle powiązane wzorem korekcyjnym.

Dalsze rozdziały poświęcimy istnieniu, jednoznaczności
i własnościom rozwiązań stochastycznych równań
różniczkowych, wzorowi Feynmana--Kaca oraz całkom
po przestrzeni trajektorii względem miary Wienera.


% ============================================================================
% ROZDZIAL 5 — SDE: konstrukcja, istnienie i jednoznacznosc
% ============================================================================
